% problem_id: S046-prop:factor_equiv_fiber
% source_id: S046 (HoTT Book)
% source_file: hlevels.tex
% statement_locator: hlevels.tex lines 1448--1467
% proof_locator: hlevels.tex lines 1469--1505
% env: lem
% pairing: tight (verified)
% status: complete (statement + author proof verbatim + reason + locators)
%
\begin{lem}\label{prop:factor_equiv_fiber}
Suppose we have a commutative diagram of functions
\begin{equation*}
  \xymatrix{
    {A} \ar[r]^{g_1} \ar[d]_{g_2} &
    {X_1} \ar[d]^{h_1} &
    \\
    {X_2} \ar[r]_{h_2}
    &
    {B}
  }
\end{equation*}
with $H:h_1\circ g_1\htpy h_2\circ g_2$, where $g_1$ and $g_2$ are $n$-connected and where $h_1$ and $h_2$ are $n$-truncated.
Then there is an equivalence
\begin{equation*}
E(H,b):\hfib{h_1}b\eqvsym\hfib{h_2}b
\end{equation*}
for any $b:B$, such that for any $a:A$ we have an identification
\[\overline{E}(H,a) :  E(H,h_1(g_1(a)))({g_1(a),\refl{h_1(g_1(a))}}) = \pairr{g_2(a),\opp{H(a)}}.\]
\end{lem}

\begin{proof}
Let $b:B$. Then we have the following equivalences:
\begin{align}
\hfib{h_1}b
& \eqvsym \sm{w:\hfib{h_1}b} \trunc n{ \hfib{g_1}{\proj1 w}}
\tag{since $g_1$ is $n$-connected}\\
& \eqvsym \Trunc n{\sm{w:\hfib{h_1}b}\hfib{g_1}{\proj1 w}}
\tag{by \cref{thm:refl-over-ntype-base}, since $h_1$ is $n$-truncated}\\
& \eqvsym \trunc n{\hfib{h_1\circ g_1}b}
\tag{by \cref{ex:unstable-octahedron}}
\end{align}
and likewise for $h_2$ and $g_2$.
Also, since we have a homotopy $H:h_1\circ g_1\htpy h_2\circ g_2$, there is an obvious equivalence $\hfib{h_1\circ g_1}b\eqvsym\hfib{h_2\circ g_2}b$.
Hence we obtain
\begin{equation*}
\hfib{h_1}b\eqvsym\hfib{h_2}b
\end{equation*}
for any $b:B$. By analyzing the underlying functions, we get the following representation of what happens to the element
$\pairr{g_1(a),\refl{h_1(g_1(a))}}$ after applying each of the equivalences of which $E$ is composed.
Some of the identifications are definitional, but others (marked with a $=$ below) are only propositional; putting them together we obtain $\overline E(H,a)$.
{\allowdisplaybreaks
\begin{align*}
\pairr{g_1(a),\refl{h_1(g_1(a))}} &
    \overset{=}{\mapsto} \Pairr{\pairr{g_1(a),\refl{h_1(g_1(a))}}, \tproj n{ \pairr{a,\refl{g_1(a)}}}}\\
  & \mapsto \tproj n { \pairr{\pairr{g_1(a),\refl{h_1(g_1(a))}}, \pairr{a,\refl{g_1(a)}} }}\\
  & \mapsto \tproj n { \pairr{a,\refl{h_1(g_1(a))}}}\\
  & \overset{=}{\mapsto} \tproj n { \pairr{a,\opp{H(a)}}}\\
  & \mapsto \tproj n { \pairr{\pairr{g_2(a),\opp{H(a)}},\pairr{a,\refl{g_2(a)}}} }\\
  & \mapsto \Pairr{\pairr{g_2(a),\opp{H(a)}}, \tproj n {\pairr{a,\refl{g_2(a)}}} }\\
  & \mapsto \pairr{g_2(a),\opp{H(a)}}
\end{align*}}
The first equality is because for general $b$, the map
\narrowequation{ \hfib{h_1}b \to \sm{w:\hfib{h_1}b} \trunc n{ \hfib{g_1}{\proj1 w}} }
inserts the center of contraction for $\trunc n{ \hfib{g_1}{\proj1 w}}$ supplied by the assumption that $g_1$ is $n$-truncated; whereas in the case in question this type has the obvious inhabitant $\tproj n{ \pairr{a,\refl{g_1(a)}}}$, which by contractibility must be equal to the center.
The second propositional equality is because the equivalence $\hfib{h_1\circ g_1}b\eqvsym\hfib{h_2\circ g_2}b$ concatenates the second components with $\opp{H(a)}$, and we have $\opp{H(a)} \ct \refl{} = \opp{H(a)}$.
The reader may check that the other equalities are definitional (assuming a reasonable solution to \cref{ex:unstable-octahedron}).
\end{proof}
